\documentclass[letterpaper,10pt]{article}

\usepackage{latexsym}
\usepackage[empty]{fullpage}
\usepackage{titlesec}
\usepackage{marvosym}
\usepackage[usenames,dvipsnames]{color}
\usepackage{verbatim}
\usepackage{enumitem}
\usepackage[hidelinks]{hyperref}
\usepackage{fancyhdr}
\usepackage[brazilian]{babel}
\usepackage{tabularx}
\input{glyphtounicode}

\usepackage[sfdefault]{roboto}

\pagestyle{fancy}
\fancyhf{}
\fancyfoot{}
\renewcommand{\headrulewidth}{0pt}
\renewcommand{\footrulewidth}{0pt}
\setlength{\footskip}{5pt}

\addtolength{\oddsidemargin}{-0.5in}
\addtolength{\evensidemargin}{-0.5in}
\addtolength{\textwidth}{1in}
\addtolength{\topmargin}{-.5in}
\addtolength{\textheight}{1.0in}

\urlstyle{same}
\raggedbottom
\raggedright
\setlength{\tabcolsep}{0in}

\titleformat{\section}{\Large\bfseries\scshape\raggedright}{}{0em}{}[\titlerule]

\pdfgentounicode=1

\newcommand{\resumeItem}[1]{\item\small{#1}}
\newcommand{\resumeSubheading}[4]{
\vspace{-1pt}\item
  \begin{tabular*}{0.97\textwidth}[t]{l@{\extracolsep{\fill}}r}
    \textbf{#1} & #2 \\
    \textit{#3} & \textit{#4} \\
  \end{tabular*}\vspace{-7pt}
}
\renewcommand\labelitemii{$\vcenter{\hbox{\tiny$\bullet$}}$}
\newcommand{\resumeSubHeadingList}{\begin{itemize}[leftmargin=0.15in, label={}]}
\newcommand{\resumeSubHeadingListEnd}{\end{itemize}}

\begin{document}

\begin{center}
  \textbf{\Huge Gabriel Frigo} \\
  \small Santo André, SP $|$ +55 11 95570-2000 $|$ \href{mailto:gabriel.frigo4@gmail.com}{gabriel.frigo4@gmail.com} $|$
  \href{https://linkedin.com/in/gabriel-frigo-b6727b275/}{linkedin.com/in/gabriel-frigo} \\
  \href{https://github.com/GabrielFrigo4}{github.com/GabrielFrigo4} $|$
  \href{https://gabrielfrigo.dev.br}{gabrielfrigo.dev.br}
\end{center}

\section{Resumo Profissional}
Estudante de Ciência da Computação na UFABC com sólida formação algorítmica, perfil autodidata e foco em \textbf{Engenharia de Software Backend e Sistemas Concorrentes}. Medalhista de Ouro na OBMEP (13º no Brasil) e competidor no ICPC e Maratona SBC de Programação. Experiência prática no desenvolvimento de APIs e microsserviços (Go, Python, SQL, C/C++), conteinerização e infraestrutura em nuvem (Docker, Linux, OCI). Praticante ativo da mentalidade \textit{AI-First}, utilizando ferramentas agênticas (\textbf{Antigravity IDE e CLI}) como aceleradores de aprendizado, validação de \textit{corner cases} e construção de software robusto.

\section{Competências Técnicas}
\resumeSubHeadingList
  \resumeItem{\textbf{Linguagens de Programação}: Go, Python, SQL, C, C++, Rust, Java, C\#, Zig, Bash, Lua}
  \resumeItem{\textbf{Backend \& Cloud}: RESTful APIs, PostgreSQL, Docker, Podman, Linux (GNU/BSD), Oracle Cloud Infrastructure (OCI), Nginx, Caddy (HTTP/2, TLS), Postman, DBeaver}
  \resumeItem{\textbf{AI-First \& Ferramentas}: Antigravity (IDE e CLI para prototipagem ágil, depuração e análise de \textit{edge cases}), Git, GitHub, GitLab, Shell Scripting}
  \resumeItem{\textbf{Engenharia \& Arquitetura}: Estruturas de Dados Avançadas, Teoria dos Grafos e Otimização Combinatória, Arquitetura de Software, Microsserviços, Concorrência e Multiprocessamento, Design Patterns, Programação Funcional e POO}
  \resumeItem{\textbf{Competências Pessoais}: Alta capacidade analítica, velocidade no domínio de novas tecnologias, questionamento crítico de soluções e comunicação colaborativa}
  \resumeItem{\textbf{Idiomas}: Português (Nativo); Inglês (Intermediário - Certificado EFSET)}
\resumeSubHeadingListEnd

\section{Formação Acadêmica}
\resumeSubHeadingList
  \resumeSubheading
    {Universidade Federal do ABC (UFABC)}{Santo André, SP}
    {Bacharelado em Ciência da Computação}{Jun/2023 -- Jun/2028}
  \resumeSubheading
    {Universidade Federal do ABC (UFABC)}{Santo André, SP}
    {Bacharelado em Ciência e Tecnologia}{Jun/2023 -- Dez/2027}
\resumeSubHeadingListEnd

\section{Experiência e Pesquisa em Engenharia}
\resumeSubHeadingList

  \resumeSubheading
    {\href{https://github.com/GabrielFrigo4/IC_Networks_Flow}{Otimização e Fluxo em Redes: Teoria e Algoritmos}}{Out/2025 -- Set/2026}
    {Iniciação Científica (CNPq / UFABC)}{Santo André, SP}
    \resumeSubHeadingList
      \resumeItem{\textbullet\ Pesquisa em Teoria dos Grafos, Algoritmos de Fluxo Máximo e Otimização Combinatória com foco em eficiência e escalabilidade.}
      \resumeItem{\textbullet\ Implementação de estruturas de dados e algoritmos em C++, com análise rigorosa de complexidade assintótica e benchmarks.}
    \resumeSubHeadingListEnd

  \resumeSubheading
    {Programação Competitiva GRUB (ICPC / Maratona SBC)}{Jan/2024 -- Presente}
    {Membro e Competidor}{UFABC, SP}
    \resumeSubHeadingList
      \resumeItem{\textbullet\ Resolução de problemas algorítmicos complexos sob limites rígidos de tempo de execução e consumo de memória utilizando C/C++.}
      \resumeItem{\textbullet\ Aplicação de programação dinâmica, grafos, geometria computacional e estruturas de dados avançadas.}
    \resumeSubHeadingListEnd

  \resumeSubheading
    {Green Team Hacker Club (GTHC)}{Dez/2023 -- Presente}
    {Analista de Projetos e Instrutor Técnico}{Santo André, SP}
    \resumeSubHeadingList
      \resumeItem{\textbullet\ Liderança técnica e mentoria para 10+ alunos sobre arquitetura de software, sistemas e boas práticas de versionamento com Git.}
      \resumeItem{\textbullet\ Ensino de C/C++, Rust e fundamentos de baixo nível aplicados à construção de software eficiente.}
    \resumeSubHeadingListEnd

  \resumeSubheading
    {\href{https://github.com/GabrielFrigo4/TT_MiniSumo_Arruela}{Robótica Competitiva Tamandutech}}{Jun/2023 -- Presente}
    {Desenvolvedor de Software}{UFABC, SP}
    \resumeSubHeadingList
      \resumeItem{\textbullet\ Desenvolvimento de lógica de software de controle de tempo real e algoritmos autônomos para tomada de decisão em robôs de competição.}
    \resumeSubHeadingListEnd

  \resumeSubheading
    {\href{https://github.com/GabrielFrigo4/PDPD_Ramsey}{Iniciação Científica PICME (IMPA / OBMEP)}}{Mar/2023 -- Jan/2026}
    {Pesquisador Acadêmico}{UFABC, SP}
    \resumeSubHeadingList
      \resumeItem{\textbullet\ Desenvolvimento de raciocínio matemático formal e combinatória aplicada à estrutura de redes e grafos.}
    \resumeSubHeadingListEnd

\resumeSubHeadingListEnd

\section{Projetos de Software}
\resumeSubHeadingList

  \resumeSubheading
    {\href{https://shop.gabrielfrigo.dev.br}{Catálogo de Lojas e Produtos (Marketplace / E-commerce)}}{Jun/2026 -- Presente}
    {Engenharia Backend \& Cloud}{Go, PocketBase, Podman, Linux, Oracle Cloud (OCI)}
    \resumeSubHeadingList
      \resumeItem{\textbullet\ Arquitetura e desenvolvimento de plataforma colaborativa de e-commerce e catálogo de produtos em produção real.}
      \resumeItem{\textbullet\ Construção de backend modular em Go com persistência de dados estruturada e isolamento de serviços via contêineres Podman.}
      \resumeItem{\textbullet\ Deploy em instância Linux na nuvem (OCI), com configuração de proxy reverso Caddy para terminação TLS/HTTPS e HTTP/2.}
      \resumeItem{\textbullet\ Aplicação ativa em produção: \href{https://shop.gabrielfrigo.dev.br/}{shop.gabrielfrigo.dev.br}}
    \resumeSubHeadingListEnd

  \resumeSubheading
    {\href{https://github.com/GabrielFrigo4/unix-sock}{Servidor Web HTTP Concorrente}}{Fev/2026 -- Mai/2026}
    {Redes, Sistemas Concorrentes e Performance}{C, POSIX/BSD Sockets, Concorrência, Caddy}
    \resumeSubHeadingList
      \resumeItem{\textbullet\ Desenvolvimento de servidor HTTP implementado do zero em C com chamadas diretas à API de POSIX Sockets.}
      \resumeItem{\textbullet\ Implementação de concorrência por multiprocessamento para tratamento simultâneo de requisições web com baixa latência.}
      \resumeItem{\textbullet\ Deploy em ambiente de produção na nuvem com Caddy proxy; Projeto live: \href{https://game.gabrielfrigo.dev.br/}{game.gabrielfrigo.dev.br}}
    \resumeSubHeadingListEnd

  \resumeSubheading
    {Mentoria em Algoritmos (Projeto Extensão OBI / UFABC)}{Jan/2026 -- Presente}
    {Instrutor de Programação}{UFABC, SP}
    \resumeSubHeadingList
      \resumeItem{\textbullet\ Treinamento de 10+ alunos em lógica, estruturas de dados e técnicas de resolução de problemas computacionais para a OBI.}
    \resumeSubHeadingListEnd

\resumeSubHeadingListEnd

\section{Prêmios e Conquistas Acadêmicas}
\resumeSubHeadingList
  \resumeItem{\textbf{Medalha de Ouro (13ª Posição Nacional, 2022)}, Prata (2018, 2019) e Bronze (2017, 2021) -- OBMEP}
  \resumeItem{\textbf{14ª Posição (2026) e 15ª Posição (2025)} -- Maratona Paulista de Programação}
  \resumeItem{\textbf{107ª Posição (2025) e 135ª Posição (2024) na Fase Regional} -- Maratona SBC de Programação / ICPC}
  \resumeItem{\textbf{Menção Honrosa (2024)} -- Olimpíada Brasileira de Informática (OBI)}
  \resumeItem{\textbf{Medalha de Prata (2023) e Menção Honrosa (2025)} -- Competição de Lógica (CELLM)}
  \resumeItem{\textbf{Medalha de Bronze (2021)} -- Olimpíada Paulista de Matemática (OPM)}
\resumeSubHeadingListEnd

\section{Certificações}
\resumeSubHeadingList
  \resumeItem{\textbullet\ \href{https://efset.org/cert/9uJjv8}{EFSET English Certificate (Intermediário)}}
\resumeSubHeadingListEnd

\end{document}
